\documentclass[11pt]{article}
\usepackage[a4paper,margin=1in]{geometry}
\usepackage{amsmath,amssymb}
\usepackage[hidelinks]{hyperref}
\title{A Disproof of Conjecture 00000008549}
\author{gaochengzhecpu (AI-assisted submission)}
\date{}
\setlength{\emergencystretch}{3em}
\begin{document}
\maketitle

The conjecture's assertion that the diamond lattice is a smallest
partition lattice with variable maximal-chain length is false.
The diamond $M_3$ is the partition lattice $\Pi_3$, and all of its
maximal chains have length $2$.  Thus it is not an example with
variable maximal-chain length at all, irrespective of minimality.

\paragraph{Precise scope.}
A partition is a collection of nonempty disjoint blocks whose union
is the underlying set.  Partitions are ordered by refinement:
$\pi\leq\sigma$ if every block of $\pi$ is contained in a block of
$\sigma$.  Equivalently, the same-block equivalence relation of $\pi$
is contained in that of $\sigma$.  We use the standard five-element
diamond $M_3$, with bottom, top, and three incomparable atoms.
The ordinary refinement order and the $\Pi_3$ diagram are also given
in the MIT lecture notes~\cite{postnikov}.
The length of a nonempty finite chain is its number of elements minus
one.  A maximal chain is maximal under inclusion among all chains.

The source's English first clause concerns graded lattices, whereas
its Chinese first clause says ordered lattices.  This proof does not
use or challenge that first clause.  It disproves the second clause,
which explicitly describes the diamond as an example with variable
length in both languages.

\paragraph{All partitions on three points.}
Use the set $\{0,1,2\}$.  A partition has either three, two, or one
block.  In the first and last cases it is forced.  In the two-block
case, one block has two elements, and there are exactly three choices
for that block.  Thus all partitions are
\[
\begin{array}{c|c|c}
\text{label}&\text{partition}&\text{rank}\\\hline
0&\{\{0\},\{1\},\{2\}\}&0\\
a&\{\{0,1\},\{2\}\}&1\\
b&\{\{0,2\},\{1\}\}&1\\
c&\{\{1,2\},\{0\}\}&1\\
1&\{\{0,1,2\}\}&2
\end{array}
\]
Each two-block partition lies above the discrete partition and below
the one-block partition.  Distinct two-block partitions do not refine
one another.  This proves $\Pi_3\cong M_3$ as ordered sets, and hence
as lattices: least upper and greatest lower bounds are determined by
the order.

\paragraph{All maximal chains.}
Let $C$ be any maximal chain of this lattice.  Both the bottom and
top must belong to $C$, since either could otherwise be added to it.
It contains at most one of $a,b,c$, since those are incomparable.
It must contain at least one of them, because $\{0,1\}$ can be
extended by inserting any atom.  Consequently the complete list of
maximal chains is
\[
0<a<1,\qquad 0<b<1,\qquad 0<c<1.
\]
Each has exactly three elements and length $2$.  Conversely, each
listed chain is maximal because its only missing elements are
incomparable with its atom.  The claimed variable-length example
therefore does not exist in the diamond.

\paragraph{Formal verification and exhaustive scope.}
The self-contained Lean 4.19.0 file imports only \texttt{Std}.
It encodes a binary relation on three points by its nine truth values
in \texttt{Fin 512}, and imposes reflexivity, symmetry, and transitivity.
The theorem \texttt{every\_prop\_partition\_encoded} shows that every
proposition-valued equivalence relation has such an encoding.
Thus the formal definition of a partition as a same-block equivalence
relation includes every actual partition.

The only equivalence-relation codes are
$273,283,341,433,511$, corresponding respectively to the five
partitions displayed above.  The maps \texttt{partitionOf} and
\texttt{label} are proved inverse.  The theorem
\texttt{arbitrary\_partition\_refinement} proves that actual relation
inclusion is precisely the diamond order.  The lattice laws for this
order and its meet and join are also verified.

All families of partitions, including arbitrary proposition-valued
families, are uniquely encoded by five membership bits.
\texttt{MaximalPartitionChain} quantifies over all such families when
expressing maximality under inclusion.  The proof transfers this
definition to the complete set of $32$ subset codes, where the only
maximal-chain codes are $19,21,25$ and each has cardinality $3$.
It also transfers maximality back and proves that these three actual
maximal chains exist, so the length conclusion is not vacuous.

The theorem \texttt{every\_maximal\_partition\_chain\_length\_two}
establishes length $2$ for every maximal chain.
The final theorem
\texttt{conjecture\_8549\_diamond\_claim\_false}
negates the existence of two actual maximal chains with unequal
lengths.  This directly negates a necessary part of the conjecture's
claim that the diamond is a variable-length example.
There are no admitted proofs or additional axioms; the final theorem
uses only \texttt{propext} and \texttt{Quot.sound}.

\begin{thebibliography}{1}
\bibitem{postnikov}
A. Postnikov, \emph{Algebraic Combinatorics, Lecture 9},
MIT 18.212, Spring 2021, section on the partition lattice.
\url{https://math.mit.edu/~apost/courses/18.212_2021/lectures/18212_lecture09.pdf}
\end{thebibliography}
\end{document}
